\documentclass[10pt]{article}
\usepackage{tmlr}
\usepackage{hyperref}
\usepackage{url}
\usepackage{graphicx}
\usepackage{booktabs}
\usepackage{multirow}
\usepackage{amsmath}
\usepackage{microtype}
\usepackage{tikz}
\usetikzlibrary{shapes.geometric, arrows.meta, positioning, fit, backgrounds}
\DeclareMathOperator*{\argmax}{arg\,max}

\title{Multi-Agent Reasoning with Consistency Verification Improves Uncertainty Calibration in Medical MCQA}

\author{\name John Ray Martinez \email jmartinez2@my.harrisburgu.edu \\
      \addr Department of Data Science and Analytics\\
      Harrisburg University of Science and Technology}

\newcommand{\fix}{\marginpar{FIX}}
\newcommand{\new}{\marginpar{NEW}}

\def\month{MM}
\def\year{YYYY}
\def\openreview{\url{https://openreview.net/forum?id=XXXX}}

\begin{document}

\maketitle

% ============================================================================
\begin{abstract}
Miscalibrated confidence scores are a practical obstacle to deploying AI
in clinical settings: a model that is always overconfident offers no useful
signal for deferral.  We present a multi-agent framework that combines
domain-specific specialist agents with Two-Phase Verification
\citep{wu2024uncertainty} and S-Score Weighted Fusion to improve both
calibration and discrimination in medical multiple-choice question answering.
Four specialist agents (respiratory, cardiology, neurology, gastroenterology)
generate independent diagnoses using Qwen2.5-7B-Instruct; each
diagnosis is then subjected to a two-phase self-verification process that
measures internal consistency and produces a Specialist Confidence Score
(S-score).  The S-scores drive a weighted fusion strategy that selects the
final answer and calibrates the reported confidence.  We evaluate across four
experimental settings: 100-question and 250-question high-disagreement subsets
of MedQA-USMLE, and 100-question and 250-question high-disagreement subsets of
MedMCQA.  On MedQA-250, our full system (Config~4) achieves 59.2\% accuracy,
ECE~$= 0.091$ (74.4\% reduction over the single-specialist baseline), and
AUROC~$= 0.630$~(+0.056).  Calibration improvements are consistent across all
four settings, including the harder MedMCQA benchmark where ECE reductions of
49--63\% are achieved even when absolute accuracy is constrained by
knowledge-intensive recall demands.  Ablation analysis across all datasets
reveals that Two-Phase Verification drives calibration improvement while
multi-agent reasoning drives accuracy gains.  Our results
establish that combining multi-agent reasoning with consistency-based
verification provides more reliable uncertainty estimates across diverse medical
question types, which is essential for safety-critical clinical AI applications.
\end{abstract}

% ============================================================================
\section{Introduction}
\label{sec:intro}

Confidence scores from medical AI systems carry real consequences: an
overconfident wrong answer may go unchallenged, while a system that hedges on
every prediction becomes clinically useless.  Yet current large language models
(LLMs) are systematically miscalibrated, producing confidence scores that
poorly track actual correctness \citep{guo2017calibration, kadavath2022language}.
On specialised, knowledge-intensive medical benchmarks this miscalibration is
especially pronounced \citep{pal2022medmcqa, jin2021disease}.

Existing recalibration methods range from post-hoc temperature scaling
\citep{guo2017calibration} and deep ensembles \citep{lakshminarayanan2017simple}
to Bayesian dropout \citep{gal2016uncertainty}, but each requires either
labelled calibration data or multiple model copies.  \citet{wu2024uncertainty}
proposed a label-free alternative: Two-Phase Verification decomposes a model's
reasoning into factual claims, answers them independently and then with
reference, and uses cross-condition inconsistency as an uncertainty signal.
Because it operates purely at inference time on the model's own outputs, it
is well-suited to deployment settings where labelled data are scarce.

Separately, multi-agent LLM systems have shown accuracy gains by pooling
diverse specialist perspectives \citep{wang2024mixture, li2024camel}, yet
how aggregation affects confidence calibration has not been studied.
Our work brings these two threads together.  We propose \textbf{MARC}
(\textbf{M}ulti-\textbf{A}gent \textbf{R}easoning with \textbf{C}onsistency
Verification), a three-component framework that combines:
\begin{enumerate}
    \item \textbf{Domain-specific specialist agents} providing diverse
          perspectives on medical questions;
    \item \textbf{Two-Phase Consistency Verification} \citep{wu2024uncertainty}
          converting per-specialist internal consistency into a Specialist
          Confidence Score (S-score);
    \item \textbf{S-Score Weighted Fusion} aggregating specialist opinions
          weighted by verified confidence to produce the final answer and
          calibrated confidence estimate.
\end{enumerate}

Our key contributions are:
\begin{itemize}
    \item \textbf{MARC}, a multi-agent architecture that systematically improves
          both calibration (ECE) and discrimination (AUROC) over single-agent
          baselines by using consistency-derived S-scores as the fusion weight,
          replacing heuristic aggregation with an uncertainty-aware signal.
    \item A rigorous four-configuration ablation study across \emph{four}
          distinct experimental settings (MedQA-100, MedQA-250, MedMCQA-100,
          MedMCQA-250), demonstrating that Two-Phase Verification is the
          primary calibration driver while multi-agent reasoning is the
          primary accuracy driver, consistently across question types and
          dataset sizes.
    \item An empirical characterisation of the knowledge-recall versus
          clinical-reasoning difficulty divide between MedMCQA and MedQA,
          and its implications for the limits of verification-based
          calibration in small LLMs.
\end{itemize}

% ============================================================================
\section{Related Work}
\label{sec:related}

\subsection{Uncertainty Quantification in Medical AI}

Uncertainty quantification (UQ) is essential in clinical AI for safe
deployment and appropriate deferral to human experts
\citep{kompa2021second, begoli2019need}.  Ensemble methods
\citep{lakshminarayanan2017simple} and Monte Carlo dropout
\citep{gal2016uncertainty} provide well-studied UQ baselines but require
multiple forward passes or model modifications.  Temperature scaling
\citep{guo2017calibration} is a simple post-hoc recalibration technique that
rescales model logits, achieving good calibration at low cost but without
improving discrimination.  Conformal prediction \citep{angelopoulos2021gentle}
offers distribution-free coverage guarantees but focuses on set-valued
predictions rather than scalar confidence scores.

\subsection{LLM Self-Evaluation and Verification}

LLMs can assess their own outputs through self-consistency sampling
\citep{wang2023selfconsistency}, chain-of-thought verification
\citep{wei2022chain}, and explicit self-critique \citep{madaan2023selfrefine}.
\citet{wu2024uncertainty} introduced Two-Phase Verification specifically for
medical QA: reasoning is decomposed into factual claims, which are then
answered independently and with reference; the cross-condition inconsistency
measures epistemic uncertainty without requiring gold labels.  We adopt this
method as our per-specialist UQ mechanism.

\subsection{Multi-Agent Systems for Medical QA}

Multi-agent LLM systems have demonstrated accuracy improvements on complex
reasoning tasks by combining diverse agent perspectives
\citep{wang2024mixture, li2024camel, liang2023encouraging}.  In the medical
domain, \citet{wang2024multi} showed that an LLM-based multi-specialist
consultation framework improves automatic diagnosis accuracy over
single-agent baselines.  \citet{du2023improving}
demonstrated that multiple agents debating their answers can reduce reasoning
errors.  However, these works focus almost exclusively on accuracy; calibration
properties of multi-agent aggregation are not systematically studied.

\subsection{Medical QA Benchmarks}

MedQA-USMLE \citep{jin2021disease} contains $\sim$12{,}700 US Medical
Licensing Examination-style questions with clinical vignettes averaging 109
words, testing integrated clinical reasoning across 4 options.  MedMCQA
\citep{pal2022medmcqa} contains over 194{,}000 questions from Indian
postgraduate medical entrance examinations (AIIMS/NEET-PG) with short stems
averaging 12.8 tokens across 21 specialties.  GPT-4 clears the USMLE passing threshold by a wide margin
\citep{nori2023capabilities} and recent frontier models exceed 90\% on
MedQA, but scores are markedly lower on MedMCQA, where specialised factual recall dominates
\citep{pal2022medmcqa, singh2025neet}.  Recent work highlights that
smaller 7B-class models exhibit pronounced accuracy drops on knowledge-intensive
recall questions, independent of reasoning capability
\citep{singh2025neet, wei2022emergent}.

% ============================================================================
\section{Methodology}
\label{sec:method}

\subsection{System Architecture}
\label{sec:arch}

Given a medical multiple-choice question with options, four specialist agents
generate independent diagnoses.  Each diagnosis undergoes Two-Phase
Verification to produce an S-score.  Finally, S-Score Weighted Fusion selects
the answer and reports a calibrated confidence.
Figure~\ref{fig:marc_architecture} illustrates the full MARC pipeline.

\begin{figure}[t]
\centering
\begin{tikzpicture}[
    font=\small,
    >=Stealth,
    every node/.style={align=center},
    io/.style={
        rectangle, rounded corners=3pt, draw, thick,
        fill=gray!15, text width=22em, minimum height=2.2em, inner sep=6pt
    },
    agent/.style={
        rectangle, rounded corners=3pt, draw, thick,
        fill=blue!10, text width=7.8em, minimum height=2em, inner sep=5pt
    },
    compbox/.style={
        rectangle, rounded corners=4pt, draw, thick,
        text width=22em, inner sep=10pt
    },
    sidelabel/.style={
        font=\footnotesize\itshape, text=gray, anchor=north west,
        align=left, text width=6.5em
    }
]

%% ── INPUT ────────────────────────────────────────────────────────────────────
\node[io] (input) {Medical Question + Answer Options};

%% ── COMPONENT 1: agents in a 2×2 grid inside a compbox ──────────────────────
\node[compbox, fill=blue!5, below=1.6em of input] (agentbox) {%
    \begin{tikzpicture}[every node/.style={align=center}, baseline]
        \node[agent] (resp) {Respiratory};
        \node[agent, right=1.2em of resp] (card) {Cardiology};
        \node[agent, below=0.5em of resp] (neur) {Neurology};
        \node[agent, right=1.2em of neur] (gast) {Gastroenterology};
    \end{tikzpicture}
};

\node[sidelabel, right=0.8em of agentbox.north east]
    {\textbf{Component 1}\\[2pt]Specialist\\Agent Team};

\draw[->] (input.south) -- (agentbox.north);

%% ── COMPONENT 2: Two-Phase Verification ─────────────────────────────────────
\node[compbox, fill=orange!10, below=1.6em of agentbox] (verif) {%
    \begin{tabular}{@{}l@{}}
        \textit{Phase 1:} specialist produces answer $a_k$, \\
        \quad reasoning $r_k$, initial confidence $C_k^{(0)}$ \\[5pt]
        \textit{Phase 2:} extract factual claims from $r_k$; \\
        \quad answer independently, then with $r_k$ in context; \\
        \quad measure cross-condition inconsistency $I_k$ \\[5pt]
        \quad $S_k \;=\; C_k^{(0)} \cdot (1 - I_k)$
    \end{tabular}
};

\node[sidelabel, right=0.8em of verif.north east]
    {\textbf{Component 2}\\[2pt]Two-Phase\\Consistency\\Verification};

\draw[->] (agentbox.south) -- (verif.north)
    node[midway, right=0.5em, font=\footnotesize\itshape, text=gray]
    {$(a_k,\,r_k,\,C_k^{(0)})$};

%% ── COMPONENT 3: S-Score Weighted Fusion ────────────────────────────────────
\node[compbox, fill=green!10, below=1.6em of verif] (fusion) {%
    \begin{tabular}{@{}l@{}}
        For each candidate $a$: \\
        \quad $\mathrm{score}(a) = |\{k : a_k\!=\!a\}| \times \bar{S}_a$ \\[5pt]
        \quad $\hat{a} = \argmax_{a}\;\mathrm{score}(a)$ \\[5pt]
        $\hat{C}$ blends vote fraction with mean and \\
        \quad min S-score of supporters
    \end{tabular}
};

\node[sidelabel, right=0.8em of fusion.north east]
    {\textbf{Component 3}\\[2pt]S-Score\\Weighted\\Fusion};

\draw[->] (verif.south) -- (fusion.north)
    node[midway, right=0.5em, font=\footnotesize\itshape, text=gray]
    {$(a_k,\;S_k)_{k=1}^{4}$};

%% ── OUTPUT ───────────────────────────────────────────────────────────────────
\node[io, below=1.6em of fusion] (output)
    {Final Answer $\hat{a}$ + Calibrated Confidence $\hat{C}$};

\draw[->] (fusion.south) -- (output.north);

\end{tikzpicture}
\caption{The MARC framework pipeline.  A medical question is distributed to
four domain-specific specialist agents (Component~1), each producing an
answer and reasoning chain.  Two-Phase Consistency Verification
(Component~2; \citealt{wu2024uncertainty}) extracts factual claims, answers
them independently and with the original reasoning in context, and measures
cross-condition inconsistency to derive an S-score $S_k$ per specialist.
S-Score Weighted Fusion (Component~3) selects the final answer by
vote-count-weighted mean S-score and derives calibrated confidence
$\hat{C}$ from the vote fraction and S-score distribution.}
\label{fig:marc_architecture}
\end{figure}

\subsection{Specialist Agents}
\label{sec:agents}

We instantiate four domain-specific specialist agents:
\begin{itemize}
    \item \textbf{Respiratory} specialist (pulmonologist)
    \item \textbf{Cardiology} specialist (cardiologist)
    \item \textbf{Neurology} specialist (neurologist)
    \item \textbf{Gastroenterology} specialist (gastroenterologist)
\end{itemize}

Each agent is driven by a Qwen2.5-7B-Instruct model \citep{qwen2025qwen25}
with a specialty-specific system prompt that instructs the agent to reason
from their domain expertise using a five-step chain-of-thought (CoT) protocol:
(1)~clinical scenario analysis, (2)~differential diagnosis, (3)~systematic
option evaluation, (4)~option comparison, and (5)~final decision with
confidence score.  Greedy decoding ($T=0$, \texttt{do\_sample=False}) is
used for all specialist calls; verification sub-calls use temperature
sampling with deterministic seeds (see Section~\ref{sec:setup}).

\subsection{Two-Phase Verification}
\label{sec:verification}

Following \citet{wu2024uncertainty}, each specialist first generates answer
$a_k$ with step-by-step explanation $r_k$ and initial confidence
$C_k^{(0)} \in [0,1]$ (Phase~1).  Phase~2 then proceeds as follows:
\begin{enumerate}
    \item \textbf{Question formulation.} Four verification questions $\{q_j\}_{j=1}^4$
          are extracted from $r_k$ by prompting the model to identify key
          factual claims.
    \item \textbf{Independent answering.} Each $q_j$ is answered without
          reference to $r_k$, yielding answers $\{a_j^{\text{ind}}\}$.
    \item \textbf{Reference answering.} Each $q_j$ is answered with $r_k$
          in context, yielding $\{a_j^{\text{ref}}\}$.
    \item \textbf{Inconsistency measurement.}  Each pair
          $(a_j^{\text{ind}}, a_j^{\text{ref}})$ is compared using
          Jaccard similarity on stop-word-filtered tokens.  A pair is
          considered consistent if the Jaccard score exceeds $\tau = 0.4$,
          or if the Jaccard score of content words (length $> 4$) exceeds
          $0.6$ (a secondary check that handles cases where one answer is
          more verbose but carries the same key terms).
          The inconsistency score over the $N_k \leq 4$ successfully
          parsed question pairs is:
          \begin{equation}
          I_k = \frac{1}{N_k}\sum_{j=1}^{N_k}
                \mathbf{1}\!\left[\text{sim}(a_j^{\text{ind}}, a_j^{\text{ref}}) < \tau\right]
          \label{eq:inconsistency}
          \end{equation}
          where $\mathbf{1}[\cdot]$ is the indicator function (1 if the
          condition holds, 0 otherwise), and $N_k = 4$ in the typical case
          when all questions are successfully parsed.
\end{enumerate}

The Specialist Confidence Score (S-score) is then computed as:
\begin{equation}
S_k = C_k^{(0)} \cdot (1 - I_k)
\label{eq:sscore}
\end{equation}
This \emph{multiplicative} formula penalises initial confidence proportionally
to measured inconsistency: a perfectly consistent specialist ($I_k=0$) retains
its full confidence, while a fully inconsistent one ($I_k=1$) receives
$S_k=0$.

\subsection{S-Score Weighted Fusion}
\label{sec:fusion}

Let $\mathcal{K}$ be the set of specialist agents, each proposing answer
$a_k \in \{A, B, C, D\}$ with S-score $S_k$.  For each candidate answer
$a$, we compute an aggregate score:
\begin{equation}
\text{score}(a) = |\{k : a_k = a\}| \;\times\; \bar{S}_a
\label{eq:fusion_score}
\end{equation}
where $\bar{S}_a = \frac{1}{|\{k:a_k=a\}|}\sum_{k:a_k=a} S_k$ is the mean
S-score of specialists supporting answer $a$, and $|\{k : a_k = a\}|$ is
their vote count.  The winning answer is $\hat{a} = \argmax_a \text{score}(a)$.

Final confidence is calibrated by the degree of consensus.  When all
specialists agree ($|\{a_k\}|=1$), we use $\hat{C} = \max_k S_k$.  Otherwise:
\begin{equation}
\hat{C} = v \cdot \bar{S}_{\hat{a}} + (1-v) \cdot \min_{k:\,a_k=\hat{a}} S_k
\label{eq:fusion_conf}
\end{equation}
where $v = |\{k : a_k = \hat{a}\}| / |\mathcal{K}|$ is the vote fraction and
$\bar{S}_{\hat{a}}$ is the mean S-score of winning supporters.  This blends
mean and minimum S-score weighted by majority strength, yielding lower
confidence for weak majorities and higher confidence for strong consensus.

\subsection{Experimental Setup}
\label{sec:setup}

\subsubsection{Datasets}

We construct four evaluation sets, all filtered to \emph{high-disagreement}
questions, specifically cases where our four specialists initially disagree, maximising
the difficulty and the importance of uncertainty quantification:

\begin{itemize}
    \item \textbf{MedQA-100}: 100 questions sampled from the high-disagreement
          subset of MedQA-USMLE \citep{jin2021disease}.  Clinical vignettes,
          mean stem length 109 words, 4 answer options.
    \item \textbf{MedQA-250}: 250 questions from the same pool.
          Serves as the primary large-scale MedQA evaluation.
    \item \textbf{MedMCQA-100}: 100 questions sampled from the high-disagreement
          subset of MedMCQA \citep{pal2022medmcqa}, spanning 21 specialties.
          Mean stem length 15 words, 4 answer options.
    \item \textbf{MedMCQA-250}: 250 questions from the same MedMCQA pool.
          Provides a larger-scale evaluation on the knowledge-intensive benchmark.
\end{itemize}

The same specialist-curation and disagreement-labelling pipeline is applied
to both benchmarks (see Appendix~\ref{app:data}).

\subsubsection{Model}

All experiments use \textbf{Qwen2.5-7B-Instruct} \citep{qwen2025qwen25} in
full FP16 precision (no quantisation).  The model is loaded with
\texttt{local\_files\_only=True} for reproducibility and runs on a single
\textbf{NVIDIA GeForce RTX~5090 Laptop GPU} (24~GB VRAM).  Greedy decoding ($T=0$,
\texttt{do\_sample=False}) is used for all specialist calls.  Temperature
sampling with deterministic seeds is used for verification sub-calls
($T_{\text{questions}}=0.3$, $T_{\text{independent}}=0.4$,
$T_{\text{reference}}=0.2$).

\subsubsection{Ablation Configurations}

We compare four configurations in a clean ablation design:
\begin{enumerate}
    \item \textbf{Config~1 -- Single Specialist (baseline)}: Respiratory
          specialist only, no verification.
    \item \textbf{Config~2 -- Single + Two-Phase}: Respiratory specialist
          with Two-Phase Verification; isolates the calibration contribution
          of verification alone.
    \item \textbf{Config~3 -- Multi + S-Score (No 2P)}: Four specialists
          fused with S-Score Weighted Fusion but without Two-Phase
          Verification (S-score $= C^{(0)}$); isolates the accuracy
          contribution of multi-agent reasoning.
    \item \textbf{Config~4 -- Multi + Two-Phase + S-Score (Full System)}:
          Four specialists with full Two-Phase Verification and S-Score
          Weighted Fusion; our main contribution.
\end{enumerate}

The progression $1\!\to\!2$ isolates Two-Phase Verification;
$1\!\to\!3$ isolates multi-agent reasoning;
$3\!\to\!4$ isolates the effect of adding verification to multi-agent.

\subsubsection{Evaluation Metrics}

\begin{itemize}
    \item \textbf{Accuracy}: Fraction of correctly answered questions.
    \item \textbf{ECE} (Expected Calibration Error, \citealt{guo2017calibration}):
          Measures alignment between predicted confidence and observed accuracy,
          using 10 equal-width bins.  Lower is better.
    \item \textbf{AUROC}: Area under the ROC curve for binary correct/incorrect
          prediction using confidence as the score.  Measures how well
          confidence discriminates correct from incorrect answers.  Higher
          is better.
\end{itemize}

% ============================================================================
\section{Results}
\label{sec:results}

\subsection{Main Results}

Table~\ref{tab:main} reports performance across all four datasets and four
configurations.  Figure~\ref{fig:accuracy_ece_auroc} visualises accuracy,
ECE, and AUROC side-by-side for each dataset.

\begin{table}[ht]
\caption{Performance comparison across all four datasets and configurations.
Best value in each column group is \textbf{bold}.
$\Delta$ columns show change relative to Config~1 baseline.
Avg Conf = mean predicted confidence score.
Time = wall-clock inference time in minutes on a single NVIDIA RTX~5090 GPU.}
\label{tab:main}
\begin{center}
\resizebox{\linewidth}{!}{%
\begin{tabular}{llcccccccc}
\toprule
\textbf{Dataset} & \textbf{Configuration} & \textbf{Acc} & $\boldsymbol{\Delta}$\textbf{Acc} & \textbf{ECE} & $\boldsymbol{\Delta}$\textbf{ECE} & \textbf{AUROC} & $\boldsymbol{\Delta}$\textbf{AUROC} & \textbf{Avg Conf} & \textbf{Time (min)} \\
\midrule
\multirow{4}{*}{MedQA-100}
  & C1: Single Specialist      & 52.0\% &  ---   & 0.374 &  ---    & 0.537 &  ---    & 0.886 &  39 \\
  & C2: Single + Two-Phase     & 52.0\% & +0.0   & 0.185 & $-$51\% & \textbf{0.678} & $+$0.141 & 0.430 &  67 \\
  & C3: Multi + S-Score (No 2P)& 55.0\% & +3.0pp & 0.356 & $-$5\%  & 0.578 & $+$0.041 & 0.896 & 152 \\
  & C4: Full System            & \textbf{59.0\%} & +7.0pp & \textbf{0.098} & $-$73\% & 0.645 & $+$0.108 & 0.549 & 249 \\
\midrule
\multirow{4}{*}{MedQA-250}
  & C1: Single Specialist      & 54.4\% &  ---   & 0.355 &  ---    & 0.574 &  ---    & 0.899 & 116 \\
  & C2: Single + Two-Phase     & 54.4\% & +0.0   & 0.178 & $-$50\% & 0.556 & $-$0.018 & 0.424 & 207 \\
  & C3: Multi + S-Score (No 2P)& 57.2\% & +2.8pp & 0.336 & $-$5\%  & 0.587 & $+$0.013 & 0.905 & 464 \\
  & C4: Full System            & \textbf{59.2\%} & +4.8pp & \textbf{0.091} & $-$74\% & \textbf{0.630} & $+$0.056 & 0.569 & 833 \\
\midrule
\multirow{4}{*}{MedMCQA-100}
  & C1: Single Specialist      & 45.0\% &  ---   & 0.469 &  ---    & 0.501 &  ---    & 0.907 &  40 \\
  & C2: Single + Two-Phase     & 45.0\% & +0.0   & 0.323 & $-$31\% & 0.437 & $-$0.064 & 0.440 &  65 \\
  & C3: Multi + S-Score (No 2P)& \textbf{52.0\%} & +7.0pp & 0.392 & $-$16\% & \textbf{0.541} & $+$0.040 & 0.902 & 181 \\
  & C4: Full System            & 50.0\% & +5.0pp & \textbf{0.237} & $-$49\% & 0.518 & $+$0.017 & 0.550 & 288 \\
\midrule
\multirow{4}{*}{MedMCQA-250}
  & C1: Single Specialist      & 42.8\% &  ---   & 0.469 &  ---    & 0.536 &  ---    & 0.897 & 102 \\
  & C2: Single + Two-Phase     & 42.8\% & +0.0   & 0.271 & $-$42\% & 0.503 & $-$0.033 & 0.460 & 167 \\
  & C3: Multi + S-Score (No 2P)& \textbf{46.8\%} & +4.0pp & 0.433 & $-$8\%  & 0.562 & $+$0.026 & 0.900 & 409 \\
  & C4: Full System            & 44.0\% & +1.2pp & \textbf{0.176} & $-$63\% & \textbf{0.594} & $+$0.058 & 0.548 & 712 \\
\bottomrule
\end{tabular}}
\end{center}
\end{table}

\begin{figure}[ht]
\begin{center}
\includegraphics[width=\linewidth]{figures/combined/combined_accuracy_ece_auroc.png}
\end{center}
\caption{Accuracy (top row), ECE (middle row), and AUROC (bottom row) for each
of the four datasets and all four configurations.  Green-outlined bars indicate
the best value in each panel.  Config~4 (Full System, dark blue) achieves the
best ECE across all four datasets, and the best accuracy on the MedQA subsets.
On MedMCQA, Config~3 achieves the best accuracy at the cost of calibration,
illustrating the knowledge-recall challenge for verification-based methods.}
\label{fig:accuracy_ece_auroc}
\end{figure}

\subsection{Calibration Analysis}

Two-Phase Verification is the dominant calibration driver.
Comparing Config~1 to Config~2 (verification only, same single agent), ECE
drops by 51\% on MedQA-100, 50\% on MedQA-250, 31\% on MedMCQA-100, and
42\% on MedMCQA-250, with no change in accuracy.  Adding verification to the
multi-agent system (Config~3$\to$4) yields even larger gains:
73\% (MedQA-100), 73\% (MedQA-250), 40\% (MedMCQA-100), and 59\% (MedMCQA-250) ECE reduction.  These improvements
are visualised in Figure~\ref{fig:calibration}.

\begin{figure}[ht]
\begin{center}
\includegraphics[width=\linewidth]{figures/combined/combined_calibration_curves.png}
\end{center}
\caption{Reliability diagrams for all four datasets.  Each panel shows
calibration curves for all four configurations against the perfect-calibration
diagonal.  Across all datasets, Config~4 (solid blue squares) lies closest to
the diagonal, confirming that Two-Phase Verification is the dominant
calibration driver regardless of question type or dataset size.}
\label{fig:calibration}
\end{figure}

\subsection{Discrimination Analysis}

Figure~\ref{fig:roc} shows ROC curves for all four datasets.  Config~4
achieves the highest AUROC on MedQA-250 (0.630, +0.056 over baseline) and
MedMCQA-250 (0.594, +0.058).  On MedQA-100, Config~2 achieves the highest
AUROC (0.678, +0.141 over baseline), with Config~4 at 0.645 (+0.108).
On MedMCQA-100, Config~3 achieves the highest AUROC (0.541) while Config~4
achieves 0.518.  The $\Delta$AUROC column in Table~\ref{tab:main} further
shows that Two-Phase Verification alone (C2) produces negative AUROC change
on three of four datasets; we analyse the mechanism behind this pattern in
Section~\ref{sec:disc_discrimination}.

\begin{figure}[ht]
\begin{center}
\includegraphics[width=\linewidth]{figures/combined/combined_roc_curves.png}
\end{center}
\caption{ROC curves for all four datasets.  Each panel plots the four
configuration curves with the random baseline (diagonal).  Config~4 achieves
the highest AUROC on MedQA-250 and MedMCQA-250.  On MedQA-100, Config~2
(single specialist with verification) achieves the highest AUROC; on
MedMCQA-100, Config~3 (multi-agent without verification) is best, illustrating
that discrimination benefits from multi-agent fusion scale more reliably with
dataset size and question type.}
\label{fig:roc}
\end{figure}

\subsection{Calibration Grid: Baseline vs.\ Full System}

Figure~\ref{fig:calgrid} provides calibration histograms comparing Config~1
(single specialist baseline) and Config~4 (full system) across all four
datasets.

\clearpage
\begin{figure}[p]
\begin{center}
\includegraphics[width=\linewidth,height=0.92\textheight,keepaspectratio]{figures/combined/combined_calibration_grid.png}
\end{center}
\caption{Calibration analysis of Qwen2.5-7B across all four evaluation sets.
Columns represent configurations C1--C4; each dataset occupies two sub-rows.
\emph{Top sub-row}: stacked confidence frequency histogram (blue = correct answer,
red = wrong answer). \emph{Bottom sub-row}: calibration histogram with 5\%-wide
bins, where bar height represents observed accuracy within that bin; the dashed
diagonal represents perfect calibration. Config~4 (Full System) consistently
produces the most evenly spread confidence distributions and aligns most closely
with the perfect-calibration diagonal across all datasets.}
\label{fig:calgrid}
\end{figure}
\clearpage

\subsection{Component Contribution Summary}

Figure~\ref{fig:component} summarises the marginal contribution of each
component (Two-Phase Verification, Multi-Agent reasoning, and their combination)
in terms of ECE reduction and AUROC change across all four datasets.

\begin{figure}[ht]
\begin{center}
\includegraphics[width=\linewidth]{figures/combined/combined_component_contribution.png}
\end{center}
\caption{(A) ECE reduction percentage and (B) AUROC change for each
component transition across all four datasets.  Adding the full system
(C1$\to$C4, dark blue) consistently delivers the largest ECE reduction
(49--74\%) across all settings.  Multi-agent reasoning (C1$\to$C3, orange)
drives the largest AUROC gains on MedQA subsets, while verification
components dominate calibration improvement throughout.}
\label{fig:component}
\end{figure}

\subsection{Ablation Summary}

Table~\ref{tab:ablation} summarises the isolated contribution of each
component on MedQA-250, the most reliable large-scale dataset.

\begin{table}[ht]
\caption{Ablation summary on MedQA-250 ($n = 250$).  Each row reports the
\emph{marginal} effect of adding one component.}
\label{tab:ablation}
\begin{center}
\small
\begin{tabular}{lccc}
\toprule
\textbf{Component Added} & $\boldsymbol{\Delta}$\textbf{Acc} & $\boldsymbol{\Delta}$\textbf{ECE} & $\boldsymbol{\Delta}$\textbf{AUROC} \\
\midrule
Two-Phase Verif.\ (C1$\to$C2) & 0.0 pp   & $-$0.177 ($-$49.9\%) & $-$0.018 \\
Multi-Agent (C1$\to$C3)       & +2.8 pp  & $-$0.019 ($-$5.4\%)  & +0.013   \\
Two-Phase added to Multi (C3$\to$C4) & +2.0 pp & $-$0.245 ($-$72.9\%) & +0.043 \\
Full System (C1$\to$C4)       & +4.8 pp  & $-$0.264 ($-$74.4\%) & +0.056   \\
\bottomrule
\end{tabular}
\end{center}
\end{table}

Multi-agent reasoning accounts for most of the accuracy gain (+2.8~pp, C1$\to$C3),
while Two-Phase Verification accounts for most of the calibration gain
($-$49.9\% ECE, C1$\to$C2; $-$72.9\% ECE, C3$\to$C4).  Together they yield
+4.8~pp accuracy and $-$74.4\% ECE on MedQA-250, with the same qualitative
pattern across all four datasets (Figure~\ref{fig:component}).

% ============================================================================
\section{Discussion}
\label{sec:discussion}

\subsection{Multi-Agent Reasoning and Accuracy}

High-disagreement questions are, by construction, cases where a single
specialist is likely to err.  Pooling four independently prompted specialists
increases the chance that at least one agent reasons from the relevant domain,
and majority voting suppresses idiosyncratic errors.  S-Score Weighted Fusion
sharpens this further: agents whose verification step reveals internal
contradictions are down-weighted, so the final answer is steered by the
specialists whose reasoning is most self-consistent.

\subsection{Two-Phase Verification and Calibration}

The unverified configurations (C1, C3) produce heavily right-skewed
confidence distributions with means near 0.90 (Figure~\ref{fig:calgrid}),
well above their actual accuracy of 43--57\%.  Two-Phase Verification brings
these means down to roughly 0.55 by penalising confidence proportionally to
measured inconsistency, making the confidence score a more faithful proxy for
correctness, as confirmed by ECE reductions of 31--74\% across all four
datasets.  This compression requires no ground-truth labels: the
signal comes entirely from the model's own reasoning.

\subsection{Calibration vs.\ Discrimination: Why Verification Alone Is Insufficient}
\label{sec:disc_discrimination}

The $\Delta$AUROC column in Table~\ref{tab:main} shows a consistent pattern:
Two-Phase Verification alone (C2) produces
\emph{negative} AUROC change on three of four datasets ($-$0.018 on
MedQA-250, $-$0.064 on MedMCQA-100, $-$0.033 on MedMCQA-250).  This
is not a contradiction with the strong ECE improvements reported for
C2.  It reflects a fundamental difference between what calibration and
discrimination measure: ECE asks whether confidence scores match observed
accuracy on average; AUROC asks whether their rank ordering correctly
separates correct from incorrect predictions.

Two-Phase Verification improves calibration by compressing all confidence
scores toward the model's actual accuracy level, reducing systematic
overconfidence.  However, this compression is proportional to measured
inconsistency, which is an imperfect proxy for incorrectness: a specialist
that produces a self-consistent but wrong answer receives a high S-score,
while one that produces an internally inconsistent but correct answer is
penalised.  The absolute values of confidence improve (lower ECE), but the
rank ordering does not reliably improve alongside them, leaving AUROC
unchanged or slightly reduced.

Multi-agent fusion restores and extends discrimination.  When four
independently prompted specialists vote and their S-scores weight the fusion,
the correct answer tends to attract specialists with higher S-scores.  This
is not because verification detects correctness directly, but because a
correct answer is more likely to elicit self-consistent reasoning across
multiple diverse agents.  The ensemble effect sharpens the separation between
confident-correct and confident-incorrect predictions, which is what AUROC
captures.  The full system (C4) consequently achieves positive AUROC gains on
all four datasets, with the largest improvements on the 250-question sets
where the ensemble signal is more stable.  Verification and multi-agent fusion
are therefore complementary rather than redundant: the former improves
calibration, the latter improves discrimination, and each is insufficient
without the other.

\subsection{MedQA vs.\ MedMCQA: Reasoning vs.\ Recall}

The full system achieves 59.2\% on MedQA-250 but only 44.0\% on MedMCQA-250
($-15.2$~pp), a gap that persists at the 100-question scale (50.0\% vs.\
52.0\% for Config~4 vs.\ Config~3 on MedMCQA-100).  Both datasets use
4-option questions with a 25\% random baseline, so the difference cannot be
attributed to the number of choices.

The structural explanation is the question type.  MedQA clinical vignettes
average 109 words and present a full patient scenario; the correct answer can
often be reached by clinical reasoning even without memorising the specific
fact.  MedMCQA stems average 15 words and test narrow factual recall from
Indian postgraduate entrance exams (AIIMS/NEET-PG), covering anatomy minutiae,
rare pharmacology, and pathology specifics across 21 specialties
\citep{pal2022medmcqa}.  At 7B parameters, the model simply does not retain
sufficient domain-specific knowledge for many of these questions
\citep{singh2025neet, kadavath2022language}.

This creates an interaction with our verification mechanism.  On MedMCQA,
moving from C3 to C4 (adding verification to multi-agent) \emph{reduces}
accuracy by 2.8~pp (250q) and 2.0~pp (100q) while still cutting ECE by 59\%
and 39\%.  The issue is that a specialist can produce internally consistent
reasoning that is nonetheless factually wrong; verification scores such
reasoning highly and may override a correct majority vote
\citep{wu2024uncertainty}.  We discuss this limitation further in
Section~\ref{sec:limitations}.

The calibration picture is more encouraging: ECE is reduced by 49--74\% in
every configuration that includes Two-Phase Verification, across all four
datasets.  Even when the model lacks the knowledge to answer correctly, it can
still be made to express appropriate uncertainty about its answer.

\subsection{Limitations}
\label{sec:limitations}

\begin{enumerate}
    \item \textbf{Consistency $\neq$ correctness.}  Two-Phase Verification
          measures internal coherence, not factual accuracy.  A specialist
          that confidently produces a self-consistent but wrong answer
          receives a high S-score.  This is the fundamental limitation of
          Wu et al.'s method and is especially problematic on knowledge-recall
          benchmarks such as MedMCQA.

    \item \textbf{Narrow domain coverage.}  Both evaluation datasets were
          pre-filtered to questions from four specialties (respiratory,
          cardiology, neurology, and gastroenterology), so results do not
          generalise to the full breadth of medical knowledge.  Domains such
          as anatomy, pharmacology, and surgery remain untested, and
          performance on those areas may differ substantially.

    \item \textbf{Computational overhead.}  Config~4 requires approximately
          $16\times$ the LLM calls of Config~1 (4 specialists, each requiring
          4 calls: 1 for the initial answer and 3 verification sub-calls for
          question formulation, independent answering, and reference answering).
          On a single NVIDIA RTX~5090 GPU, MedQA-250 takes 833~minutes for
          Config~4 vs.\ 116~minutes for Config~1 (7.2$\times$ wall-clock
          overhead); MedMCQA-250 takes 712~minutes vs.\ 102~minutes
          (7.0$\times$).  Parallel specialist inference is an obvious optimisation path.
\end{enumerate}

% ============================================================================
\section{Future Work}
\label{sec:future}

The most consequential open problem is the one identified in
Section~\ref{sec:limitations}: a specialist can produce self-consistent but
factually wrong reasoning, and the verification step has no way to catch it.
Retrieval augmentation is the most direct path forward.  Grounding Phase~2
verification against a medical knowledge base such as PubMed would replace
pure internal consistency with evidence-anchored consistency, allowing the
S-score to penalise claims that contradict established medical knowledge
rather than only claims that contradict the specialist's own reasoning.

Two questions about the current design remain empirically open.  First,
whether the calibration improvements hold at larger model scales (13B, 70B,
GPT-4 class): larger models are better calibrated to begin with
\citep{kadavath2022language}, so the absolute ECE gains may shrink even if
the relative benefit of verification persists.  Second, whether the
vote-count fusion can be improved with learned weights: a small specialist
performance validation set could provide signal for per-specialty reliability
that the current heuristic ignores.  Both questions are testable with the
infrastructure already in place.

% ============================================================================
\section{Conclusion}
\label{sec:conclusion}

We studied whether combining domain-specific specialist agents with
consistency-based self-verification improves uncertainty calibration in
medical multiple-choice QA.  Across four evaluation sets drawn from MedQA
and MedMCQA, the full system (Config~4) reduces ECE by 49--74\% over the
single-specialist baseline and improves AUROC by up to +0.108.  The
ablation confirms that multi-agent reasoning drives most of the accuracy
improvement while Two-Phase Verification drives most of the calibration
improvement.

The MedMCQA results tell a different story.  Even when absolute accuracy
stagnates because the 7B model lacks the specialised factual knowledge that
MedMCQA demands, the calibration mechanism continues to function: the system
produces lower confidence on the answers it gets wrong.  Accuracy and
calibration can therefore be improved independently, and consistency-based
verification remains a practical deferral signal even in knowledge-constrained
settings.

The central open problem is the confidence penalty the system assigns to
self-consistent but factually wrong answers.  Grounding the verification step
in external medical knowledge is the most direct path to resolving it.

% ============================================================================
\subsubsection*{Broader Impact Statement}

The system studied here is a research prototype and should not be used for
actual clinical decisions without extensive prospective validation.
The accuracy levels achieved (44--59\%) are well below what autonomous
clinical decision support would require.  Better-calibrated confidence scores
could help clinicians identify cases where AI recommendations warrant
additional scrutiny, but this benefit depends on the calibration holding
in deployment settings not covered by our evaluation.  The specialist agents'
knowledge is bounded by the pretraining data of a 7B model, which may
reflect geographic, demographic, and linguistic biases in that corpus.

\subsubsection*{Acknowledgments}

This work is part of a doctoral dissertation at Harrisburg University of
Science and Technology.  The author thanks Dr.~Vaida for guidance and support.

\subsubsection*{Data and Code Availability}

Evaluation datasets (MedQA-USMLE, MedMCQA) are publicly available.
Experimental code and result files will be released upon publication at
\url{https://github.com/[anonymised]}.

% ============================================================================
\bibliography{main}
\bibliographystyle{tmlr}

% ============================================================================
\appendix

\section{Appendix}

\subsection{Dataset Construction}
\label{app:data}

High-disagreement subsets are constructed by the following three-step
procedure applied identically to both benchmarks:

\begin{enumerate}
    \item \textbf{Specialist curation.}  A lightweight curation pass runs all
          four specialist agents on the full source datasets (MedQA training
          split: $\sim$10{,}000 questions; MedMCQA training split:
          $\sim$49{,}000 questions).  Each agent is prompted to respond with
          only the option letter (A/B/C/D) and no explanation
          (\texttt{max\_new\_tokens=8}); only the letter is needed to detect
          disagreement, so this imposes no quality loss relative to the main
          experiments.  Early stopping halts after collecting 220
          high-disagreement and 60 agreement candidates.
    \item \textbf{Disagreement labelling.}  A question is labelled
          high-disagreement if at least two of the four specialists propose
          different answers.
    \item \textbf{Subset construction.}  From the curated pool, 100-question
          and 250-question subsets are drawn by taking the top-ranked
          high-disagreement questions (sorted by number of differing specialist
          answers) and a random sample of agreement questions (fixed seed for
          reproducibility), retaining only questions with valid ground-truth
          labels (\texttt{cop} $\in \{1,2,3,4\}$ for MedMCQA;
          \texttt{answer\_idx} present for MedQA).  The combined set is
          shuffled before evaluation.
\end{enumerate}

\subsection{Prompt Templates}
\label{app:prompts}

\subsubsection*{Specialist Agent — System Prompt}

\begin{verbatim}
You are an expert medical specialist in {specialty}.
You have deep knowledge in this specific domain and are
consulting on a medical case.

{knowledge_context}

Your task is to:
1. Analyze the question from your specialty's perspective
2. Provide your expert opinion on the correct answer
3. Explain your reasoning clearly
4. Include relevant medical knowledge from your specialty

Be precise, evidence-based, and focus on your area of expertise.
\end{verbatim}

\subsubsection*{Specialist Agent — User Prompt}

\begin{verbatim}
Question: {question}

Options:
{options}

As a {specialty} specialist, please use CHAIN-OF-THOUGHT reasoning:

STEP 1: Understand the clinical scenario
- What are the key symptoms, signs, or findings?
- What is the patient's presentation?
- What is the clinical context?

STEP 2: Consider differential diagnoses from your specialty
- What conditions in your specialty could explain this?
- What are the diagnostic criteria for each?
- What are the distinguishing features?

STEP 3: Evaluate each option systematically
- For EACH option, evaluate:
  * Is this medically correct?
  * Does this fit the clinical scenario?
  * What evidence supports or refutes this?

STEP 4: Compare options
- Which option best fits the clinical scenario?
- Which option has the strongest evidence?

STEP 5: Make your decision
- Select the most appropriate answer
- Provide confidence score (0-1)

IMPORTANT: For ANSWER, provide ONLY the option LETTER
(A, B, C, or D). Do NOT provide the full text.

Format your response as:
STEP_1_ANALYSIS: [clinical scenario analysis]
STEP_2_DIFFERENTIAL: [differential diagnoses]
STEP_3_OPTION_EVALUATION: [evaluation of each option]
STEP_4_COMPARISON: [comparison of options]
STEP_5_DECISION: [final decision reasoning]
ANSWER: [Single letter: A, B, C, or D]
CONFIDENCE: [0.0-1.0]
REASONING: [detailed explanation]
\end{verbatim}

\subsubsection*{Two-Phase Verification — Question Formulation Prompt}

\begin{verbatim}
You are a medical verification expert. Your task is to
formulate verification questions based on the explanation.

Question: {question}
Proposed Answer: {answer}
Explanation: {reasoning}

Based on the explanation above, formulate exactly 4 specific
verification questions that check the factual claims made in
the explanation. These questions should:
1. Target specific medical facts or claims in the explanation
2. Be answerable independently (without the explanation)
3. Help verify the correctness of the reasoning
4. Cover the key factual claims in the explanation

Format your response as:
VERIFICATION_QUESTIONS:
1. [First question]
2. [Second question]
3. [Third question]
4. [Fourth question]
\end{verbatim}

\subsubsection*{Two-Phase Verification — Independent Answering Prompt}

\begin{verbatim}
You are a medical expert. Answer the following verification
questions based on your medical knowledge.

Original Question Context: {question}

Verification Questions:
{verification_questions}

IMPORTANT: Answer these questions based on your medical
knowledge ONLY. Do NOT reference any specific explanation
or reasoning from the original question.

Format your response as:
ANSWERS:
1. [Answer to first question]
2. [Answer to second question]
3. [Answer to third question]
4. [Answer to fourth question]
\end{verbatim}

\subsubsection*{Two-Phase Verification — Reference Answering Prompt}

\begin{verbatim}
You are a medical expert. Answer the following verification
questions while referencing the provided explanation.

Original Question Context: {question}

Original Explanation:
{reasoning}

Verification Questions:
{verification_questions}

IMPORTANT: Answer these questions while referencing the
explanation above. Use the explanation to guide your answers.

Format your response as:
ANSWERS:
1. [Answer to first question]
2. [Answer to second question]
3. [Answer to third question]
4. [Answer to fourth question]
\end{verbatim}

\subsection{S-Score Formula Variants}
\label{app:sscore}

Three S-score formula variants were considered:
\begin{align}
S^{\text{mult}} &= C^{(0)} \cdot (1 - I) \label{eq:mult} \\
S^{\text{wavg}} &= \alpha \, C^{(0)} + (1-\alpha)(1-I),\quad \alpha=0.65 \label{eq:wavg}\\
S^{\text{pure}} &= 1 - I \label{eq:pure}
\end{align}
The multiplicative formula (Eq.~\ref{eq:mult}) was selected because it
preserves the sign relationship between initial confidence and the final
S-score: a high-confidence specialist is penalised more in absolute terms
than a low-confidence one for the same level of inconsistency.  The weighted
average (Eq.~\ref{eq:wavg}) decouples the two signals additively and requires
tuning $\alpha$; the pure consistency formula (Eq.~\ref{eq:pure}) discards
initial confidence entirely, which performed poorly on questions where
the initial confidence was already well-calibrated.

\end{document}
